\chapter{1906: Henri Moissan}
\label{chap:chemistry1906}

The Nobel Prize in Chemistry for the year 1906 was awarded to Henri Moissan.

\textbf{Motivation:}
in recognition of the great services rendered by him in his investigation and isolation of the element fluorine, and for the adoption in the service of science of the electric furnace called after him

\section{Henri Moissan}

\subsection*{Early Life and Education}

Henri Moissan was born on 1852-09-28 in Paris, France.
Ferdinand Frédéric Henri Moissan was a French chemist and pharmacist who won the 1906 Nobel Prize in Chemistry for his work in isolating fluorine from its compounds. Among his other contributions, Moissan discovered moissanite and contributed to the development of the electric arc furnace. Moissan was one of the original members of the International Atomic Weights Committee.

Moissan came from a family of modest means and was obliged to earn his living while studying. He was apprenticed to a watchmaker and later to a pharmacist, and it was through pharmacy that he first encountered chemistry. He studied under Edmond Frémy at the Muséum National d'Histoire Naturelle and at the Sorbonne, where he was appointed professor in 1900. The combination of his practical apprenticeship and his academic training gave him the technical skill that his later work demanded.

\subsection*{Scientific Context}

The isolation of fluorine, the most reactive of all elements, was one of the most difficult problems in chemistry. Humphry Davy, Michael Faraday, and many others had attempted it and failed, because fluorine reacts violently with almost every material that can contain it, including the vessels of the apparatus itself. The problem required a complete command of the chemistry of fluorine compounds and, above all, an experimental skill that could prevent the apparatus from being destroyed as fast as it was built. The successful isolation of fluorine in 1886 was therefore regarded as the solution of a problem that had defied the best chemists of Europe for more than seventy years.

\subsection*{Career and Major Research}

Moissan was professor of inorganic chemistry at the Sorbonne. In 1886 he succeeded in isolating fluorine, the most reactive of the halogens, by electrolysis of anhydrous hydrogen fluoride, and he studied the reactions of this element. He also invented the electric arc furnace, which enabled the synthesis of refractory materials such as carbides and artificial diamonds.

The invention of the electric arc furnace in 1892 opened up a new region of temperature for chemistry. By passing a powerful electric current through a carbon arc enclosed in a furnace of lime, Moissan obtained temperatures of more than three thousand degrees, and in this furnace he melted and reduced refractory oxides that had resisted all previous attempts. He prepared the carbides of many metals, including calcium carbide, whose hydrolysis produces acetylene, and he studied the behavior of the metals themselves at the highest temperatures. In 1893 he announced the production of small artificial diamonds by dissolving carbon in molten iron and cooling it rapidly under great pressure, a result that, although its interpretation was later questioned, excited enormous public interest.

\subsection*{Nobel Prize-winning Work}

The work for which Henri Moissan was awarded the Nobel Prize in Chemistry in 1906 was described by the Nobel Committee as: "in recognition of the great services rendered by him in his investigation and isolation of the element fluorine, and for the adoption in the service of science of the electric furnace called after him".

This research fundamentally advanced the field by isolating elemental fluorine and developing the electric furnace for high-temperature chemistry, enabling the synthesis of artificial diamonds and new refractory materials.

The isolation of fluorine was achieved by the electrolysis of a solution of potassium hydrogen fluoride in anhydrous hydrogen fluoride, carried out in an apparatus of platinum with electrodes also of platinum and iridium. The choice of materials was crucial: platinum withstands the attack of fluorine better than almost any other metal, and the low temperature maintained in the cell kept the corrosive power of the gas within bounds. On June 26, 1886, Moissan collected the first sample of fluorine gas, a discovery announced to the French Academy of Sciences within a week.

\subsection*{The Work in Detail}

The experimental problem in isolating fluorine lay not only in preparing the gas but in showing that it was fluorine. Moissan collected it in a carefully dried platinum apparatus rather than over water, which would react with the gas. He identified it by its chemical action on silicon, which it attacks immediately, and by its ability to set fire to sulfur and phosphorus. He then determined its most important properties: it is a pale yellow-green gas, slightly denser than air, and it combines directly with nearly every element. The difficulty and danger of the work were extreme, and Moissan is said to have suffered from the effects of fluorine poisoning in later years.

In the years that followed, Moissan explored the chemistry of fluorine with great thoroughness. He showed that fluorine attacks gold and platinum at higher temperatures, that it combines directly with most metals and nonmetals, and that it displaces the other halogens from their compounds. He prepared many new fluorine compounds, including fluorides of sulfur and iridium, and he worked out practical methods for handling the gas. This program established the chemistry of fluorine as a coherent branch of the science.

\subsection*{Reception}

The isolation of fluorine was greeted as a triumph of experimental chemistry, and Moissan received the prize of the Academy of Sciences, the Lacaze Prize, within a few months of his discovery. His work with the electric furnace extended his reputation into the new field of high-temperature chemistry, and he was elected to the French Academy of Sciences in 1891. The Nobel Prize of 1906, awarded the year before his death, was the recognition of a career in which a notorious experimental problem and a new technique of high-temperature investigation had both been brought to success.

\subsection*{Significance and Impact}

The impact of Henri Moissan's work extends far beyond the specific achievement recognized by the Nobel Prize.
The isolation of fluorine, the most reactive element, was a triumph of experimental technique. Henri Moissan's electric furnace enabled high-temperature synthesis and the production of artificial diamonds, opening new avenues in materials science.

The electric furnace transformed the technology of the chemical industry as well as the science. Moissan's high-temperature work contributed to the development of calcium-carbide manufacture and the production of acetylene for lighting and welding, while the furnace made possible the study and production of many metals, carbides, and alloys that could not otherwise be obtained. The combination of high temperature with the use of carbon as both electrode and reducing agent established methods that remain in use in modern metallurgical and abrasives industries.

\subsection*{Later Career and Honors}

Henri Moissan continued his scientific work until 1907-02-20, passing away in Paris, France.
He received numerous honors including membership in major scientific academies and societies.
Henri Moissan was elected to the French Academy of Sciences and received numerous international honors.

Moissan's later years were devoted largely to the electric furnace and to the study of the metals and carbides prepared in it. He received the Davy Medal of the Royal Society in 1896 and was created a commander of the Legion of Honor; his laboratory in Paris attracted students from many countries. He died in 1907, shortly after the award of the Nobel Prize, while still in full possession of his scientific powers.

\subsection*{Legacy}

The legacy of Henri Moissan endures through the lasting impact of his scientific contributions.
The Moissan process for fluorine production and the electric furnace technique enabled modern high-temperature chemistry and materials science. The Moissan crater on the Moon bears Henri Moissan's name.

Moissan's isolation of fluorine is still regarded as one of the great landmarks of experimental inorganic chemistry, and the element he freed has become one of the most industrially important in the modern world. The electric furnace, developed in his laboratory, became the standard instrument of high-temperature chemistry and metallurgy, and the carbide chemistry he founded led directly to modern applications from acetylene welding to the production of abrasives and semiconductors.

\subsection*{Collaborators and Students}

Moissan trained a school of French chemists who carried on his methods in inorganic and high-temperature chemistry, and his laboratory was among the most active in Europe in the study of fluorine compounds and refractory materials. His work on the atomic weights of fluorine and other elements made him a valued member of the International Atomic Weights Committee, and his careful determinations contributed to the accuracy of the periodic table in its early years.

\subsection*{Modern Developments}

Moissan's isolation of fluorine led to the modern fluorine chemistry industry, including Teflon, hydrofluorocarbons, refrigerants, and many fluorinated pharmaceuticals. His invention of the electric arc furnace enabled high-temperature chemistry and the production of carbides, including silicon carbide, a key modern abrasive and semiconductor material.

The chemistry of fluorine that Moissan founded has grown into a major branch of modern industry. Fluoropolymers line cookware and insulate cables, fluorinated gases serve as refrigerants and in electronics manufacture, and fluorinated compounds are prominent in pharmaceutical and agricultural chemistry. The electric arc furnace has developed into the industrial arc furnaces that produce steel, ferroalloys, and silicon, and the mineral moissanite, the natural form of silicon carbide first identified in his name, has become a well-known gemstone produced synthetically. Each of these applications descends from the experimental achievements of Moissan's laboratory.

\subsection*{Selected Publications}

"Le Four Électrique" (1897)
"Sur le diamant artificiel" (1893)
"Le Fluor et ses Composés" (1900)
